\documentclass[10pt,a4paper]{amsart}
\usepackage[margin=19mm]{geometry}
\usepackage{amsmath,amssymb,mathtools}
\usepackage[scr=rsfs,cal=euler]{mathalfa}
\usepackage{xfrac}
\usepackage[unicode,hidelinks,bookmarks=false]{hyperref}
\newcommand{\ie}{i.e.\ }
\newcommand{\cf}{cf.\ }
\newcommand{\resp}{respectively\ }
\newcommand{\Subsection}{\subsection}
\providecommand{\nobreakdash}{\nobreak\hbox{-}}
\setlength{\emergencystretch}{2em}
\multlinegap0pt
\usepackage{bm}
\usepackage{bbm}
\usepackage{dsfont}
\usepackage{xspace}
\usepackage{cancel}
\usepackage{mathtools}
\usepackage{amsthm}
\makeatletter
\@ifundefined{theorem}{\newtheorem{theorem}{Theorem}}{}
\@ifundefined{proposition}{\newtheorem{proposition}[theorem]{Proposition}}{}
\makeatother
\newcommand{\mc}[1]{\mathcal{#1}}
\title[Selecting Interpretable Circular Coordinates from Data]{Selecting Interpretable Circular Coordinates from Data}
\author{Vincent P. Grande, Marina Meila}
\date{Source: arXiv:2607.08230v1 (CC BY 4.0)}
\begin{document}
\maketitle

\noindent\emph{Source and licence.} Selecting Interpretable Circular Coordinates from Data. Version arXiv:2607.08230v1 (CC BY 4.0), distributed under the Creative Commons Attribution 4.0 International licence, as recorded in the arXiv licence metadata for this exact version and shown on the abstract page https://arxiv.org/abs/2607.08230v1. Full rights notice: licenses/arxiv-2607.08230-CC-BY-4.0.txt.\par\smallskip

\noindent\textit{Editorial scope.} Counted unit: the Proposition whose source label is \texttt{None} in the article (source0.tex lines 1238--1254, source label \texttt{None}), delivered together with the article's own complete original proof of it (source0.tex lines 1255--1516, one \texttt{proof} environment of 262 lines). No mathematics is written, repaired or re-derived: statement and proof are the article's own text line for line. Required context retained uncounted: eq:kernel-mass-expansion (lines 1060--1063); eq:def-sampled-cochains (lines 1177--1182); eq:discrete-inner-product-expanded (lines 1193--1200). Every definition, display and label of those lines is retained unchanged. Omitted from this package and not redistributed: 1--1059, 1064--1176, 1183--1192, 1201--1237, 1517--1635. Nothing inside a retained excerpt is omitted or altered: every retained body is delivered exactly as the article's lines are written, so no omission receipt of any kind accompanies this package. Citations: the retained excerpts carry no citation command at all, so no bibliography entry of the article is transcribed and no reference is redistributed. Reference adaptation: none was needed and none was made; every reference inside the delivered unit resolves inside it, so no unresolved reference or citation is produced. Printed slips kept literal: no printed word, symbol, hypothesis or number of the counted statement or of its proof was altered. Layout and class: the article's own class is \texttt{scrartcl} with packages including biblatex, fontenc, amsmath, amssymb, mathtools, mathpazo, relsize, amsthm, enumerate, algorithm, algpseudocode, csquotes, verbatim, graphicx; the wrapper is the lane's pinned \texttt{amsart} editorial wrapper at 10pt on A4, which restates the theorem-like environments the retained text needs and adds the article's own macro definitions that the retained text needs (\texttt{\textbackslash mc}), with extra wrapper packages (bm, bbm, dsfont, xspace, cancel, mathtools). The delivered numbering is the unit's own. Assets: no retained span contains a figure environment, a graphics-inclusion command, a PDF-page inclusion, a table or a TikZ picture; no image file is copied into this package, the package declares no asset dependency and no image file is redistributed with it. Licence: version arXiv:2607.08230v1 of this article is distributed under the Creative Commons Attribution 4.0 International licence, as recorded in the arXiv licence metadata for this exact version and shown on the abstract page https://arxiv.org/abs/2607.08230v1. A full rights notice is delivered at licenses/arxiv-2607.08230-CC-BY-4.0.txt. Characters: every retained excerpt is pure ASCII, so the delivered engine, XeLaTeX with its default Latin Modern text font, reports no missing character.
\begin{equation}
\label{eq:kernel-mass-expansion}
q_h(x)=m_0\pi(x)+O(h^2)
\end{equation}

\begin{equation}
\label{eq:def-sampled-cochains}
c_\alpha(e)\coloneq \alpha(x_i,x_j),
\qquad
c_\beta(e)\coloneq \beta(x_i,x_j).
\end{equation}

\begin{equation}
\label{eq:discrete-inner-product-expanded}
\langle c_\alpha,c_\beta\rangle_M
=
\frac{2dm_0^2}{m_2h^2}
\sum_{\substack{e=(x_i,x_j)\in E_1\\ K_h(x_i,x_j)>0}}
\frac{K_h(x_i,x_j)}{q_iq_j}\,c_\alpha(e)c_\beta(e).
\end{equation}

\begin{proposition}[Sampling error]
Assume that
\begin{enumerate}
\item $(\mc{M},g)$ is a smooth compact $d$-dimensional Riemannian manifold without boundary;
\item the sample $X=\{x_1,\dots,x_N\}$ is drawn i.i.d. from a measure $\nu=\pi\,\mu$, where $\pi\in C^2(\mc{M})$ is strictly positive;
\item the kernel profile $\kappa$ is bounded, supported in $[0,1)$, and satisfies $m_2>0$;
\item $h=h_N\to 0$, $N h_N^d/\log N\to \infty$, and $h_N<\mathrm{inj}(\mc{M})$ for all sufficiently large $N$;
\item the $1$-skeleton of the chosen simplicial complex on $X$ contains every unordered pair $\{x_i,x_j\}$ for which $K_{h_N}(x_i,x_j)\neq 0$.
\end{enumerate}
Let $M=M_{N,h_N}$ be the diagonal matrix defined in Step~4 from this sample and bandwidth $h_N$.
Then, for every fixed smooth $\alpha,\beta\in \Omega^1(\mc{M})$, the sampled cochains from \eqref{eq:def-sampled-cochains} satisfy
\begin{equation}
\langle c_\alpha,c_\beta\rangle_M-I_{h_N}(\alpha,\beta)
=
O_{\mathbb P}\!\left(\sqrt{\frac{\log N}{N h_N^d}}\right).
\end{equation}
\end{proposition}

\begin{proof}
For convenience, write
\begin{equation}
r_i\coloneq
\begin{cases}
\dfrac{N}{q_i}, & \text{if } q_i>0,\\[0.8ex]
0, & \text{if } q_i=0,
\end{cases}
\qquad
s_i\coloneq \frac{1}{q_h(x_i)},
\qquad
\Delta_{N,h}\coloneq \sup_{1\le i\le N}|r_i-s_i|.
\end{equation}
The proof has two steps.

\paragraph{Step 1: Replace the empirical reciprocals}
Fix $i$. Conditional on $x_i$, the random variables
\begin{equation}
Y_{i\ell}\coloneq K_h(x_i,x_\ell),
\qquad
\ell\neq i,
\end{equation}
are i.i.d.\ by assumption \textup{(2)}, bounded by $\|\kappa\|_\infty h^{-d}$ by assumption \textup{(3)}, and satisfy
\begin{equation}
\mathbb E[Y_{i\ell}\mid x_i]=q_h(x_i),
\qquad
\mathbb E[Y_{i\ell}^2\mid x_i]\le \|\kappa\|_\infty h^{-d}q_h(x_i)\le Ch^{-d}.
\end{equation}
Hence Bernstein's inequality yields constants $c,C>0$, independent of $i$, such that for every $0<t\le 1$,
\begin{equation}
\label{eq:conditional-bernstein}
\mathbb P\!\left(
\left|
\frac{1}{N-1}\sum_{\ell\neq i}Y_{i\ell}-q_h(x_i)
\right|>t \,\middle|\, x_i
\right)
\le
2\exp\!\left(-cNh^dt^2\right).
\end{equation}
Taking expectations and applying a union bound over $i$ gives
\begin{equation}
\label{eq:q-over-n-minus-one}
\sup_{1\le i\le N}
\left|
\frac{q_i}{N-1}-q_h(x_i)
\right|
=
O_{\mathbb P}\!\left(\sqrt{\frac{\log N}{N h^d}}\right).
\end{equation}
Since $q_i/N=\frac{N-1}{N}\cdot q_i/(N-1)$ and $q_h$ is uniformly bounded, the difference between $q_i/N$ and $q_i/(N-1)$ is $O_{\mathbb P}(1/N)$ uniformly in $i$, hence
\begin{equation}
\label{eq:q-over-n}
\sup_{1\le i\le N}
\left|
\frac{q_i}{N}-q_h(x_i)
\right|
=
O_{\mathbb P}\!\left(\sqrt{\frac{\log N}{N h^d}}\right).
\end{equation}

By \eqref{eq:kernel-mass-expansion}, whose proof uses the geometric hypotheses in assumption \textup{(1)}, and by the lower bound on $\pi$ from assumption \textup{(2)}, there exists $c_0>0$ such that $q_h(x)\ge c_0$ for all $x$ and all sufficiently small $h$. Since assumption \textup{(4)} gives $\sqrt{\log N/(N h^d)}\to 0$, \eqref{eq:q-over-n} implies that the event
\begin{equation}
\sup_i\left|\frac{q_i}{N}-q_h(x_i)\right|\le \frac{c_0}{2},
\end{equation}
has probability tending to $1$. On this event we also have $q_i/N\ge c_0/2$ for every $i$. Therefore,
\begin{equation}
\left|r_i-s_i\right|
=
\left|
\frac{1}{q_i/N}-\frac{1}{q_h(x_i)}
\right|
\le
\frac{2}{c_0^2}
\left|
\frac{q_i}{N}-q_h(x_i)
\right|,
\end{equation}
and \eqref{eq:q-over-n} implies
\begin{equation}
\label{eq:reciprocal-rate}
\Delta_{N,h}
=
O_{\mathbb P}\!\left(\sqrt{\frac{\log N}{N h^d}}\right).
\end{equation}

Define
\begin{equation}
\label{eq:def-I-tilde}
\widetilde I_{N,h}(c_\alpha,c_\beta)
\coloneq
\frac{2dm_0^2}{m_2h^2}
\sum_{e=(x_i,x_j)\in E_1}
\frac{K_h(x_i,x_j)}{N^2q_h(x_i)q_h(x_j)}c_\alpha(e)c_\beta(e).
\end{equation}
This is well-defined because assumption \textup{(3)} gives $m_2>0$.
Because $r_i$ and $s_i$ are uniformly $O_{\mathbb P}(1)$, the product difference satisfies
\begin{equation}
\sup_{i,j}|r_ir_j-s_is_j|
=
O_{\mathbb P}(\Delta_{N,h}).
\end{equation}
Substituting this bound into \eqref{eq:discrete-inner-product-expanded} and \eqref{eq:def-I-tilde} yields
\begin{align}
\label{eq:difference-I-Itilde}
\left|\langle c_\alpha,c_\beta\rangle_M-\widetilde I_{N,h}(c_\alpha,c_\beta)\right|
&\le
C\Delta_{N,h}\,
\frac{1}{N^2h^2}
\sum_{e=(x_i,x_j)\in E_1}
K_h(x_i,x_j)\,|c_\alpha(e)c_\beta(e)|.
\end{align}
Because assumption \textup{(3)} says that $\kappa$ is supported in $[0,1)$ and assumption \textup{(4)} gives $h<\mathrm{inj}(\mc{M})$ for all sufficiently large $N$, every contributing pair satisfies $d_{\mc{M}}(x_i,x_j)<h$ and the geodesic integrals defining $c_\alpha(e)$ and $c_\beta(e)$ are well-defined. Since $\mc{M}$ is compact by assumption \textup{(1)}, smoothness of the fixed forms $\alpha,\beta$ gives the uniform bound
\begin{equation}
\label{eq:cochain-size}
|c_\alpha(e)|+|c_\beta(e)|\le Ch.
\end{equation}
Therefore
\begin{equation}
\frac{1}{N^2h^2}
\sum_{e=(x_i,x_j)\in E_1} K_h(x_i,x_j)\,|c_\alpha(e)c_\beta(e)|
\le
\frac{C}{N^2}
\sum_{1\le i\neq j\le N} K_h(x_i,x_j).
\end{equation}
The expectation of this quantity is
\begin{equation}
\frac{1}{N^2}\sum_{i\neq j}\mathbb E[K_h(X_i,X_j)]
=
\frac{N-1}{N}\,\mathbb E[q_h(X_1)]
\le C,
\end{equation}
so Markov's inequality gives
\begin{equation}
\frac{1}{N^2h^2}
\sum_{e\in E_1}K_h(x_i,x_j)\,|c_\alpha(e)c_\beta(e)|
=
O_{\mathbb P}(1).
\end{equation}
Together with \eqref{eq:reciprocal-rate}, this proves
\begin{equation}
\label{eq:reciprocal-replacement-error}
\langle c_\alpha,c_\beta\rangle_M-\widetilde I_{N,h}(c_\alpha,c_\beta)
=
O_{\mathbb P}\!\left(\sqrt{\frac{\log N}{N h^d}}\right).
\end{equation}

\paragraph{Step 2: Compare the pairwise average with the continuum expectation}
After the reciprocal replacement, the discrete sum becomes a pairwise average of the kernel
\begin{equation}
\label{eq:def-Phi-h}
\Phi_h(x,y)
\coloneq
\frac{dm_0^2}{m_2h^2}
\frac{K_h(x,y)}{q_h(x)q_h(y)}
\alpha(x,y)\beta(x,y).
\end{equation}
\begin{equation}
\label{eq:def-U-N-h}
U_{N,h}
\coloneq
\frac{2}{N(N-1)}
\sum_{1\le i<j\le N}\Phi_h(x_i,x_j).
\end{equation}
Because $\alpha(y,x)=-\alpha(x,y)$ and $\beta(y,x)=-\beta(x,y)$, the product $\alpha(x,y)\beta(x,y)$ is symmetric in $(x,y)$, hence so is $\Phi_h$. Assumption \textup{(5)} therefore implies
\begin{equation}
\label{eq:Itilde-vs-U}
\widetilde I_{N,h}(c_\alpha,c_\beta)
=
\frac{N-1}{N}\,U_{N,h}.
\end{equation}

On the support of $K_h$, smoothness gives $|\alpha(x,y)|+|\beta(x,y)|\le Ch$, hence
\begin{equation}
\label{eq:phi-moments}
\|\Phi_h\|_\infty\le C h^{-d},
\qquad
\mathbb E[\Phi_h(X_1,X_2)^2]\le C h^{-d}.
\end{equation}
Here the support condition in assumption \textup{(3)} and the injectivity-radius condition in assumption \textup{(4)} ensure that $\alpha(x,y)$ and $\beta(x,y)$ are evaluated only on geodesic pairs with $d_{\mc{M}}(x,y)<h$, and compactness from assumption \textup{(1)} gives the uniform constants. The second bound follows from the first because assumption \textup{(3)} also gives $K_h^2\le \|\kappa\|_\infty h^{-d}K_h$.

Set
\begin{equation}
g_h(x)\coloneq \mathbb E[\Phi_h(x,X_2)]-I_h(\alpha,\beta),
\qquad
H_h(x,y)\coloneq \Phi_h(x,y)-I_h(\alpha,\beta)-g_h(x)-g_h(y).
\end{equation}
Then $\mathbb E[g_h(X_1)]=0$ and
\begin{equation}
\mathbb E[H_h(x,X_2)]=\mathbb E[H_h(X_1,y)]=0.
\end{equation}
Because the sample points are i.i.d.\ by assumption \textup{(2)}, the Hoeffding decomposition gives
\begin{equation}
\label{eq:hoeffding}
U_{N,h}-I_h(\alpha,\beta)
=
\frac{2}{N}\sum_{i=1}^N g_h(x_i)
+
\frac{2}{N(N-1)}
\sum_{1\le i<j\le N} H_h(x_i,x_j).
\end{equation}
Moreover,
\begin{equation}
|g_h(x)|
\le
\int_{\mc{M}}|\Phi_h(x,y)|\,d\nu(y)+|I_h(\alpha,\beta)|
\le C,
\end{equation}
because $|\alpha(x,y)\beta(x,y)|\le Ch^2$ on the support of $K_h$ and $\int K_h(x,y)\,d\nu(y)=q_h(x)=O(1)$ uniformly. Hence
\begin{equation}
\label{eq:linear-part-variance}
\mathbb E\!\left[
\left(
\frac{2}{N}\sum_{i=1}^N g_h(X_i)
\right)^2
\right]
=
O\!\left(\frac{1}{N}\right).
\end{equation}

For the canonical part, mixed second moments vanish unless the two unordered pairs coincide: if $\{i,j\}\neq \{k,\ell\}$, then either the pairs are disjoint, in which case independence gives zero, or they share exactly one index, in which case conditioning on the shared variable and using the canonical property again gives zero. Therefore
\begin{equation}
\label{eq:canonical-second-moment}
\mathbb E\!\left[
\left(
\frac{2}{N(N-1)}
\sum_{1\le i<j\le N} H_h(X_i,X_j)
\right)^2
\right]
\le
\frac{C}{N^2}\,\mathbb E[\Phi_h(X_1,X_2)^2]
=
O\!\left(\frac{1}{N^2h^d}\right).
\end{equation}
By Chebyshev's inequality, \eqref{eq:linear-part-variance}, and \eqref{eq:canonical-second-moment},
\begin{equation}
\label{eq:u-stat-rate}
U_{N,h}-I_h(\alpha,\beta)
=
O_{\mathbb P}\!\left(\frac{1}{\sqrt{N}}\right)
+
O_{\mathbb P}\!\left(\frac{1}{N h^{d/2}}\right)
=
O_{\mathbb P}\!\left(\sqrt{\frac{1}{N h^d}}\right).
\end{equation}
The last equality uses $h\to 0$, so $h^d\le 1$ for all sufficiently large $N$ and therefore $N^{-1/2}\le (Nh^d)^{-1/2}$.
%This direct second-moment estimate is enough here; more general exponential inequalities for order-two U-statistics are available in Houdr\'e and Reynaud-Bouret \cite[Theorem~3.4]{HoudreReynaudBouret2003}.

Since \eqref{eq:Itilde-vs-U} implies
\begin{equation}
\widetilde I_{N,h}(c_\alpha,c_\beta)-U_{N,h}
=
O_{\mathbb P}\!\left(\frac{1}{N}\right),
\end{equation}
because $U_{N,h}=I_h(\alpha,\beta)+O_{\mathbb P}((Nh^d)^{-1/2})=O_{\mathbb P}(1)$, we obtain
\begin{equation}
\label{eq:Itilde-minus-Ih}
\widetilde I_{N,h}(c_\alpha,c_\beta)-I_h(\alpha,\beta)
=
O_{\mathbb P}\!\left(\sqrt{\frac{1}{N h^d}}\right).
\end{equation}
Combining \eqref{eq:reciprocal-replacement-error} and \eqref{eq:Itilde-minus-Ih} proves the claim.
\end{proof}

\end{document}
